% content_pool.tex — example seed items showing the metadata-header convention.
%
% Each item is a single \newcommand whose body is one or more
% \cventry / \subcventry / \cvitem calls. The LaTeX comment block above each
% item is metadata the /write-cv skill reads (without parsing TeX):
%
%   %%% ITEM: <macro-name-without-backslash>
%   %%% always: <true|false>
%   %%% group: <string>          (optional — omit when not applicable)
%   %%% relevance: <jobtype>=<high|medium|low>, ...
%
% Section (Berufserfahrung / Ausbildung / Projekte) is derived from the
% nearest preceding  % ===== <name> =====  block header.
%
% This example pool is intentionally small (2 jobs, 2 degrees, 3 projects) to
% keep the demo compile to ~3 pages. The real pool lives in
% ~/application-pipeline/data/user-info/content_pool.tex and grows freely.

% ===== Berufserfahrung =====

%%% ITEM: itemJobExample1
%%% always: false
%%% relevance: games=high, mle=medium
\newcommand{\itemJobExample1}{%
  \cventry[0.5em]{<<JOB_1_DATES>>}{<<JOB_1_ROLE>>}{<<JOB_1_EMPLOYER>>}{}{}{%
    \begin{itemize}
      \item <<JOB_1_BULLET_1>>
      \item <<JOB_1_BULLET_2>>
      \item <<JOB_1_BULLET_3>>
    \end{itemize}}%
}

%%% ITEM: itemJobExample2
%%% always: false
%%% relevance: games=high, mle=low
\newcommand{\itemJobExample2}{%
  \cventry{<<JOB_2_DATES>>}{<<JOB_2_ROLE>>}{<<JOB_2_EMPLOYER>>}{}{}{%
    \begin{itemize}
      \item <<JOB_2_BULLET_1>>
      \item <<JOB_2_BULLET_2>>
    \end{itemize}}%
}

% ===== Ausbildung =====

%%% ITEM: itemDegreeMaster
%%% always: true
%%% relevance: games=high, mle=high
\newcommand{\itemDegreeMaster}{%
  \cventry{<<DEGREE_1_DATES>>}{<<DEGREE_1_TITLE>>}{<<DEGREE_1_SCHOOL>>}{}{}{}%
  \subcventry{<<DEGREE_1_THESIS_LABEL>>}{<<DEGREE_1_THESIS_TITLE>>}{<<DEGREE_1_THESIS_DESC>>}%
  \vspace{0.5em}%
}

%%% ITEM: itemDegreeBachelor
%%% always: true
%%% relevance: games=medium, mle=high
\newcommand{\itemDegreeBachelor}{%
  \cventry{<<DEGREE_2_DATES>>}{<<DEGREE_2_TITLE>>}{<<DEGREE_2_SCHOOL>>}{}{}{}%
  \subcventry{<<DEGREE_2_THESIS_LABEL>>}{<<DEGREE_2_THESIS_TITLE>>}{<<DEGREE_2_THESIS_DESC>>%
    \vspace{-.8em}%
    \begin{itemize}\item <<DEGREE_2_AWARD>>\end{itemize}}%
}

% ===== Projekte =====

%%% ITEM: itemProjectMLE
%%% always: false
%%% relevance: mle=high, games=low
\newcommand{\itemProjectMLE}{%
  \cventry[0.5em]{<<PROJECT_MLE_DATE>>}{<<PROJECT_MLE_TITLE>>}{}{}{}{<<PROJECT_MLE_DESC>>}%
}

%%% ITEM: itemProjectGames
%%% always: false
%%% relevance: games=high, mle=low
\newcommand{\itemProjectGames}{%
  \cventry[0.5em]{<<PROJECT_GAMES_DATE>>}{<<PROJECT_GAMES_TITLE>>}{}{}{}{<<PROJECT_GAMES_DESC>>}%
}

%%% ITEM: itemProjectAwarded
%%% always: true
%%% relevance: games=high, mle=medium
\newcommand{\itemProjectAwarded}{%
  \cventry[0.5em]{<<PROJECT_AWARDED_DATE>>}{<<PROJECT_AWARDED_TITLE>>}{}{}{}{<<PROJECT_AWARDED_DESC>>%
    \vspace{-.5em}%
    \begin{itemize}\item <<PROJECT_AWARDED_NOTE>>\end{itemize}}%
}
